Theorem 5. No arc separates $\mathbf{R}^{2}$.\\
Proof. Let $A$ be an $\operatorname{arc}$ in $\mathbf{R}^{2}$. Since $A$ is bounded, $\mathbf{R}^{2}-A$ has exactly one unbounded component. Thus we need to show that $\mathbf{R}^{2}-A$ has no bounded component.

If $U$ is a bounded component of $\mathbf{R}^{2}-A$, then $\operatorname{Fr} U \subset A$, and $\operatorname{Fr} U$ is closed. On an interval $[a, b]$, every closed set $M$ lies on a minimal interval\\
$\left[a^{\prime}, b^{\prime}\right]$, with $a^{\prime}, b^{\prime} \in M$. (The reason is that the least upper bound and the greatest lower bound of $M$ must belong to $M$; these are $b^{\prime}$ and $a^{\prime}$.) It follows that every homeomorphic image of $[a, b]$ has the same property; $A$ has a subarc $A^{\prime}$ which contains $\operatorname{Fr} U$ and has its end-points in $\operatorname{Fr} U$. We may therefore assume, with no loss of generality, that the end-points $T, T^{\prime}$ of $A$ lie in $\mathrm{Fr} U$.

We now enclose $A$, as usual, in a 2-cell $\bar{I}$, such that $A$ intersects $\operatorname{Fr} I$ in exactly two points $P, R$. As in Figure 4.8, these may not be the end-points

\begin{center}
\fbox{\parbox{0.82\linewidth}{\centering \emph{Figure 4.8 omitted in this rendered excerpt; see the raw book source above.}}}
\end{center}


$T, T^{\prime}$ of $A$. By the preceding theorem, there is a broken line $B$, from $S$ to $Q$, such that $B \cap \operatorname{Fr} I=\{Q, S\}$, Int $B \subset I, B \cap A_{1}=\varnothing$, and $B \cap A_{3}=\varnothing$.

Let $V$ be the first point of $B$ (in the order from $S$ ) that lies in $A_{2}$; let $W$ be the last point of $B$ that lies in $A_{2}$; let $B_{1}$ be the arc from $S$ to $V$ in $B$; let $B_{2}$ be the arc from $V$ to $W$ in $A_{2}$; let $B_{3}$ be the arc from $W$ to $Q$ in $B$; and let $B^{\prime}=B_{1} \cup B_{2} \cup B_{3}$. By Theorem 2, $P$ and $R$ lie in the boundaries of different components of $I-B^{\prime}$. Therefore $T$ and $T^{\prime}$ lie in different components of $I-B^{\prime}$. But this is impossible, because $U$ is connected, $U \cap B^{\prime}=\varnothing$, and $T$ and $T^{\prime}$ are limit points of $U$. This contradiction completes the proof of the theorem. \(\square\)
